\documentclass[10pt,a4paper]{article}
\usepackage{amsmath,amssymb,amsthm,graphicx}
\usepackage{makecell}
\usepackage{dsfont}
% Card-local minimal preamble (journal unit sigma-2025-023-macdonald-identities-weyl-kac-denominator-fo, P2-ahl16 precedent).
%% Verbatim from the pinned source: the author theorem environments and macro block,
%% so every macro and environment of the counted spans is available.
% Omitted and why: the documentclass line (the journal class sigma is not on the host;
%% the assembler supplies the class) and page-geometry/font blocks (they fight the
%% forced xelatex engine).
\usepackage{amsmath,amssymb,amsthm,mathtools}

\renewcommand{\textfraction}{0.01}
\renewcommand{\topfraction}{0.99}
\renewcommand\theadalign{bc}
\renewcommand\theadfont{\bfseries}
\renewcommand\theadgape{\Gape[4pt]}
\renewcommand\cellgape{\Gape[4pt]}
\newtheorem{Theorem}{Theorem}[section]
\newtheorem{Corollary}[Theorem]{Corollary}
\newtheorem{Lemma}[Theorem]{Lemma}
\newtheorem{Proposition}[Theorem]{Proposition}
 { \theoremstyle{definition}
\newtheorem{Definition}[Theorem]{Definition}
\newtheorem{Note}[Theorem]{Note}
\newtheorem{Example}[Theorem]{Example}
\newtheorem{Remark}[Theorem]{Remark} }
\numberwithin{equation}{section}
\newcommand{\ti}{\tilde}
\newcommand{\al}{\alpha}
\newcommand{\om}{\omega}
\newcommand{\ga}{\gamma}
\newcommand{\la}{\lambda}
\newcommand{\eps}{\varepsilon}
\newcommand{\si}{\sigma}
\newcommand{\La}{\Lambda}
\newcommand{\Ga}{\Gamma}
\newcommand{\ov}[1]{\overline{#1}}
\newcommand{\os}[2]{\overset{#1}{#2}}
\newcommand{\scb}[2]{\scalebox{#1}{#2}}
\newcommand{\crossp}{\overset{+}{\leadsto}}
\newcommand{\crossn}{\overset{-}{\leadsto}}
\newcommand{\fh}{\mathfrak{h}}
\newcommand{\fg}{\mathfrak{g}}
%
\newcommand{\cC}{\mathcal{C}}
\newcommand{\cF}{\mathcal{F}}
\newcommand{\cU}{\mathcal{U}}
\newcommand{\cM}{\mathcal{M}}
\newcommand{\cA}{\mathcal{A}}
\newcommand{\cB}{\mathcal{B}}
\def\cch{\mathcal H}
\def\ccsc{\mathcal {SC}}
\def\ccdd{\mathcal {DD}}
\def\ccd{\mathcal {D}}
\def\ccp{\mathcal P}
\def\ccc{\mathcal C}
\DeclareMathOperator{\BGP}{BG}
\DeclareMathOperator{\WLP}{WL}
\DeclareMathOperator{\core}{core}
\DeclareMathOperator{\quot}{quot}
\newcommand{\scal}[2]{\langle #1,#2\rangle}
\newcommand{\scald}[2]{(#1,#2)}
\newcommand{\qu}[1]{\qquad\text{#1}\qquad}
\newcommand{\te}[1]{\text{ #1 }}
\newcommand{\Z}{\mathbb{Z}}
\newcommand{\N}{\mathbb{N}}
\newcommand{\C}{\mathbb{C}}
\newcommand{\R}{\mathbb{R}}
\newcommand{\K}{\mathbb{K}}
\newcommand{\Q}{\mathbb{Q}}
\newcommand{\bB}{\mathbb{B}}
\newcommand{\bA}{\mathbb{A}}
\newcommand{\fB}{\mathfrak{B}}
\newcommand{\wB}{\widehat{B}}
\newcommand{\mk}{\mathrm{k}}
\newcommand{\mK}{\mathrm{K}}
\newcommand{\sLa}{\mathsf{\La}}
\newcommand{\cAJ}{\mathcal{A}_J}
\DeclareMathOperator{\Sp}{Sp}
\DeclareMathOperator{\GL}{GL}
\DeclareMathOperator{\OO}{SO}
\renewcommand{\OE}{\mathrm{O}}
\DeclareMathOperator{\I}{I}
\DeclareMathOperator{\oo}{so}
\renewcommand{\oe}{\mathrm{o}^{\text{even}}}
\newcommand{\func}[1]{\operatorname{#1}}
\newcommand{\lra}{\longrightarrow}
\newcommand{\cev}[1]{\reflectbox{\ensuremath{\vec{\reflectbox{\ensuremath{#1}}}}}}
\newcommand{\cI}{\mathcal{I}}
\newcommand\blfootnote[1]{%
	\begingroup
	\renewcommand\thefootnote{}\footnote{#1}%
	\addtocounter{footnote}{-1}%
	\endgroup
}
\newcommand{\syc}{\mathsf{s_C}}
\newcommand{\lb}{\llbracket}
\newcommand{\rb}{\rrbracket}
\usepackage{cleveref}
\crefname{originallemma}{original lemma}{original lemmas}
\Crefname{originallemma}{Original lemma}{Original lemmas}
\makeatletter
\newlabel{lem:indices}{{18.2}{source}{}{sourcepointer-1}{}}
\newlabel{lem:indices@cref}{{[originallemma][0][]18.2}{[1][0][]source}{}{}{}}
\makeatother
\crefname{originalequation}{original equation}{original equations}
\Crefname{originalequation}{Original equation}{Original equations}
\makeatletter
\newlabel{eq:ineg}{{7.2}{source}{}{sourcepointer-2}{}}
\newlabel{eq:ineg@cref}{{[originalequation][0][]7.2}{[1][0][]source}{}{}{}}
\makeatother
\title{Inversion sets of affine Grassmannian elements and hook lengths of the corresponding n-core}
\author{Lecouvey, C\textbackslash{}'\{e\}dric, Wahiche, David}
\date{doi:10.3842/SIGMA.2025.023; CC BY 4.0}
\begin{document}
\maketitle
\setcounter{section}{7}
\setcounter{Theorem}{8}
\setcounter{equation}{2}
\noindent\textit{Editorial note.} All delivered mathematical text is the author's own, copied from the pinned journal source file sources/sigma-2025-023-macdonald-identities-weyl-kac-denominator-fo/pinned-original-source.tex (Macdonald Identities, Weyl-Kac Denominator Formulas and Affine Grassmannian Elements); no statement or proof step was rewritten, reordered or removed.\par\medskip
\noindent\textit{Editorial note.} The counted proof is the author's complete argument: The proof invokes two numbered named inputs of the same article, declared as external navigation pointers: Lemma lem:indices (18.2, invoked at line 3111) and equation eq:ineg ((7.2), invoked at line 3116). The map J\_n is defined by the unnumbered display at lines 3078-3083 immediately preceding the statement (standing setup, not a numbered object). The source writes Psi\textasciicircum{}\{-1\} (capital) in the statement and proof but psi\textasciicircum{}\{-1\} (lowercase, Definition def:psi 18.1) at line 3085; carried verbatim as written. The commented-out orange author note at lines 3101-3103 is carried verbatim and does not render. The census printed field 32.3 is the published-version coordinate; the card face renders Proposition 7.9 by source-text counter stepping. No numbered equations are rendered inside the span.\par\medskip
\noindent\textit{Editorial note.} The label prop:inversion names the counted statement, and the proof environment immediately following it in the source is delivered verbatim as its proof.\par\medskip
\noindent\textit{Editorial note.} The counted statement is carried by the author's own Proposition environment in the article source.\par\medskip
\noindent\textit{Editorial note.} Journal-unit assembly: card-local preamble (author theorem environments and macros verbatim), article class, no bibliography in the counted span (the span carries no \textbackslash{}cite commands).\par\medskip
\noindent\textit{Editorial note.} Revision r002: the r001 card-local preamble was malformed - the lane's preamble extractor dropped five structural lines of the author's \textbackslash{}blfootnote definition (\textbackslash{}begingroup, \textbackslash{}addtocounter, \textbackslash{}endgroup, the closing brace, a bare \% comment), leaving an unbalanced brace and an xelatex failure; the extractor is fixed and the preamble is regenerated verbatim from the pinned source. Additionally \textbackslash{}usepackage\{makecell\} is declared via assembly.extra\_packages because the source preamble redefines the journal's tablehead macros (\textbackslash{}theadalign, \textbackslash{}theadfont, \textbackslash{}theadgape, \textbackslash{}cellgape via \textbackslash{}Gape) that the article class does not define; no card-body text changed.\par\medskip
\noindent\textit{Editorial note.} Revision r003: the r002 assembly failed on the undefined macro \textbackslash{}mathds (used verbatim in the author proof at source line 3112, the indicator function); \textbackslash{}usepackage\{dsfont\} is declared via assembly.extra\_packages (the journal class provides it; the article class does not). No card-body text changed.\par\medskip
\noindent\textit{Editorial scope.} the counted statement and its own proof environment only\par\medskip
\noindent\textit{Allowed standard prerequisites (editorial).}\begin{itemize}
\item The map J\_n: the affine Grassmannian elements of type A\_\{n-1\}\textasciicircum{}(1) to binary words, defined by the unnumbered display at source lines 3078-3083 (c\_k = 1 iff k >= w(i)-1 with k = w(i)-1 mod n), immediately preceding the statement.
\item Lemma lem:indices (printed 18.2): the index correspondence between boxes of the core partition and congruence classes of the binary word.
\item Equation eq:ineg (printed (7.2)): the inequality system closing the count, established just before the statement.
\item The inversion set R(w) of an affine permutation and its height grading ht\_n.
\end{itemize}\medskip
\ifdefined\hypertarget\hypertarget{sourcepointer-1}{}\fi
\noindent\textit{Source reference.} Original lemma 18.2: sources/sigma-2025-023-macdonald-identities-weyl-kac-denominator-fo/pinned-original-source.tex, source lines 1829--1829. Navigation pointer invoked at line(s) 3111: the counted text invokes that lemma of the same article (the index correspondence between word positions and boxes) as a named input, which is not retained in this card, so the reference is declared as a pointer carrying the number the article itself prints for it (the section and the place of the statement inside that section, derived by the census v3.5 rule for the shared section-scoped theorem counter); the cited lemma is an allowed input and none of it is reproduced inside the span.\par\medskip
\ifdefined\hypertarget\hypertarget{sourcepointer-2}{}\fi
\noindent\textit{Source reference.} Original equation 7.2: sources/sigma-2025-023-macdonald-identities-weyl-kac-denominator-fo/pinned-original-source.tex, source lines 3073--3073. Navigation pointer invoked at line(s) 3116: the counted text cites that numbered equation of the same article, which is not retained in this card, so the reference is declared as a pointer carrying the number the article itself prints for it (the section and the place of the equation inside that section, derived by the census v3.5 rule for the section-scoped equation counter); the cited equation is an allowed input and none of it is reproduced inside the span.\par\medskip
\begin{Proposition}\label{prop:inversion}
Set $w\in\widetilde{\mathfrak{S}}^0_n$ and $\lambda=\Psi^{-1}(J_n(w))$ the corresponding $n$-core partition. Define the set $R_i(w):=\bigl\lbrace \alpha\in R(w),\, ht_n\bigl(w^{-1}(\alpha)\bigr)=n-i\bigr\rbrace$.
Then for any $i\in\lbrace 1,\dots,n-1\rbrace$, we have
\[\#R_i(w)=\#\lbrace h\in \cch(\lambda),\, h=n-i\rbrace.\]
%$$R_n(w)=\lbrace \alpha\in R(w), ht_n(w^{-1}(\alpha))=n-i\rbrace=
\end{Proposition}

\begin{proof}
Set $w\in\widetilde{\mathfrak{S}}^0_n$, $J_n(w)=(c_k)_{k\in\Z}$ and $\lambda=\Psi^{-1}(J_n(w))$ the corresponding $n$-core partition. Set $\alpha\in R_i(w)$ and $(a,b)\in\lbrace 1,\dots,n\rbrace^2$ such that $\alpha=\varepsilon_a-\varepsilon_b+k\delta$
and $w^{-1}(a)-w^{-1}(b)\equiv n-i\pmod{n}$. Take $(c,d)\in\lbrace 1,\dots,n\rbrace$ the unique elements such that $a\equiv w(c)\pmod{n}$ and $b\equiv w(d)\pmod{n}$. Take $(m_a,m_b)\in\Z^2$ such that $a=w(c)+m_an$ and $b=w(d)+m_bn$. Thus
\[\alpha=\varepsilon_{w(c)}-\varepsilon_{w(d)}+(k+m_b-m_a)\delta.\]
%Therefore
%$\textcolor{orange}{ht_n(w^{-1}(\alpha))= n-i}$.
%Moreover, we have that $\textcolor{orange}{w^{-1}}(a)=c+m_an$ and $\textcolor{orange}{w^{-1}}(b)=d+m_bn$.
Hence, by definition of $J_n$, we have
\begin{align*}
&\lbrace k\in \Z,\, c_{k}=1,\, k\equiv w(c)-1\pmod{n}\rbrace=\lbrace a-1+qn,\, q\geq -m_a\rbrace,\\
&\lbrace k\in \Z,\, c_{k}=0,\, k\equiv w(d)-1\pmod{n}\rbrace=\lbrace b-1+qn,\, q< -m_b\rbrace.
\end{align*}


By Lemma \ref{lem:indices}, we conclude the proof by noticing that the set of boxes $s\in\lambda$ such that $h_s=n-i$ such that the indices $i_s$ and $j_s$, are congruent to $w(c)-1$ respectively $w(d)-1$ $\pmod{n}$ corresponds to the following set:
\[\lbrace (a-1+kn,b-1+(k+\mathds{1}_{a>b}(a,b))n),\, -m_a\leq k \leq -m_b-1-\mathds{1}_{a>b}(a,b)\rbrace.
\]


Therefore, setting $k''=k+m_a+\mathds{1}_{a>b}(a,b)$, $k''$ follows the inequalities \eqref{eq:ineg}, which concludes the proof.
\end{proof}

\end{document}
